\documentclass[11pt]{article}
\usepackage[margin=1in]{geometry}
\usepackage{amsmath,amssymb}
\usepackage[hidelinks]{hyperref}

\newcommand{\Sha}{\operatorname{Sha}}
\title{Disproof of Conjecture 00000000255\\
\large An elliptic curve with negative root number and a square-order obstruction}
\author{}
\date{}

\begin{document}
\maketitle

Consider the elliptic curve
\[
  E/\mathbb Q:\quad y^2+y=x^3-x.
\]
Its global root number is \(-1\). If its Tate--Shafarevich group
is finite, the Cassels--Tate pairing forces its order to be a
positive square. The two assertions about \(\Sha(E)\) in the
conjecture therefore cannot hold together.

\section*{The root number}

The integral Weierstrass coefficients are
\((a_1,a_2,a_3,a_4,a_6)=(0,0,1,-1,0)\), giving
\((b_2,b_4,b_6,b_8)=(0,-2,1,-1)\). In particular
\[
  \Delta=-8(-2)^3-27=37,\qquad c_4=-24(-2)=48.
\]
This model is minimal at \(37\) because the discriminant valuation
is one and changes between Weierstrass models occur in multiples
of twelve. It has good reduction at every other prime.
At \(37\), completing the square with \(Y=2y+1\) gives
\[
  Y^2=4x^3-4x+1=4(x-5)^2(x+10)\quad\text{over }\mathbb F_{37}.
\]
The double root \(5\) is distinct from the simple root \(-10\).
Writing \(u=x-5\), the tangent cone of the node at \((5,0)\)
is \(Y^2-23u^2\), since \(4(5+10)\equiv23\pmod{37}\).
The residue \(23\) is not a square. In fact successive squaring
gives \(23^2\equiv11\), \(23^4\equiv10\),
\(23^8\equiv26\), and \(23^{16}\equiv10\pmod{37}\),
so \(23^{18}\equiv-1\pmod{37}\). A nonzero square would have
eighteenth power \(1\) by Fermat's little theorem.
Thus the two tangents are distinct but are not defined over
\(\mathbb F_{37}\), giving nonsplit multiplicative reduction.

The good and nonsplit multiplicative finite places have local
root number \(+1\). The real place has local root number \(-1\).
Multiplying the local signs gives \(w(E/\mathbb Q)=-1\)
\cite[Example~3.22 and Section~8.1]{dokchitser}.

\section*{The order constraint on \(\Sha(E)\)}

For an elliptic curve, the Cassels--Tate pairing is alternating
and is nondegenerate after quotienting by the maximal divisible
subgroup \cite[Introduction]{poonen-stoll}. A finite group has
zero divisible subgroup. Thus, if \(G=\Sha(E)\) is finite,
the pairing on \(G\) is perfect and alternating.

For completeness, a finite abelian group with such a pairing
has square order. Choose \(x\in G\) of maximal order \(n\).
If \(n=1\) the assertion is immediate. Otherwise perfection
implies that the character \(\langle x,-\rangle\) has order \(n\).
Its image in \(\mathbb Q/\mathbb Z\) is cyclic of order \(n\),
so one can choose \(y\) with \(\langle x,y\rangle\) of order
\(n\). Then \(y\) also has order \(n\), by maximality of \(n\).
Pairing a relation between \(x\) and \(y\) with each of them
shows that both coefficients are divisible by \(n\). Consequently
\(H=\langle x,y\rangle\) has order \(n^2\), and its pairing is
nondegenerate. Every character on \(H\) is represented by an element
of \(H\), so each element of \(G\) differs from an element of \(H\)
by an element of \(H^\perp\). Hence \(G=H\oplus H^\perp\), and the
pairing on \(H^\perp\) is again perfect and alternating. Induction
on group order proves the square-order assertion.

An equality \(a^2=2b^2>0\) would force both \(a\) and \(b\) to
be even, producing a smaller positive solution after division.
Thus if \(\Sha(E)\) is finite, its positive square order cannot
be twice a square as asserted for root number \(-1\).
If \(\Sha(E)\) is infinite, the finiteness assertion already fails.

\section*{Lean verification}

The Lean~4.19.0 project checks the Weierstrass invariants, the
factorization and tangent computation modulo \(37\), and the
product of local signs under the stated reduction rules. It
proves by descent that a positive square cannot be twice a square
and derives the contradiction from the finite Cassels--Tate
square-order theorem. Run \verb|lake build| in \verb|lean/| to
verify the formal results.

\begin{thebibliography}{2}
\bibitem{dokchitser}
T.~Dokchitser, \emph{Notes on the parity conjecture},
\href{https://arxiv.org/abs/1009.5389}{arXiv:1009.5389}.
\bibitem{poonen-stoll}
B.~Poonen and M.~Stoll,
\emph{The Cassels--Tate pairing on polarized abelian varieties},
Annals of Mathematics \textbf{150} (1999), 1109--1149.
\href{https://math.mit.edu/~poonen/papers/sha.pdf}{Author's copy}.
\end{thebibliography}

\end{document}
